% A display inside a vertical box, and what starts a paragraph.
%
% `tex.web` §1090 starts a paragraph from *vertical* mode. Math mode is not
% vertical mode, and a character in a formula is a `\mathchar`, never the thing
% that begins a paragraph.
%
% That distinction is invisible on the page and decisive inside a `\vbox`.
% Opening a display breaks the paragraph so far (§1145), which leaves the box
% with no paragraph in progress — and we asked the *box* whether a paragraph
% was needed rather than the mode. So the first character of the formula
% started one: `\everypar` fired inside the display, `\parskip` went in, and the
% display's own `\predisplaypenalty`, `\abovedisplayskip`, equation box,
% `\postdisplaypenalty` and `\belowdisplayskip` were all appended to that
% paragraph's horizontal list — which §1206 `resume_after_display` clears.
%
% Every display inside a vertical box vanished, contents and all. On the page it
% worked, because there the mode really is `mmode` and the question never came
% up, so no gate built on whole documents could see it.
\catcode`\{=1 \catcode`\}=2 \catcode`\#=6 \catcode`\$=3
\catcode`\^=7 \catcode`\_=8
\font\rip=cmr10 \font\mi=cmmi10 \font\sy=cmsy10 \font\ex=cmex10
\textfont0=\rip \scriptfont0=\rip \scriptscriptfont0=\rip
\textfont1=\mi \scriptfont1=\mi \scriptscriptfont1=\mi
\textfont2=\sy \scriptfont2=\sy \scriptscriptfont2=\sy
\textfont3=\ex \scriptfont3=\ex \scriptscriptfont3=\ex
\rip
\hsize=200pt \parindent=0pt \baselineskip=12pt \lineskip=1pt \lineskiplimit=0pt

% A paragraph with a display in the middle of it: two text lines and the
% display between them.
\setbox0\vbox{\noindent A$$A$$A\par}
\message{<para-with-display=\the\ht0>}

% The same paragraph with no display, for the difference to be readable.
\setbox1\vbox{\noindent AA\par}
\message{<para-alone=\the\ht1>}

% A display that ends the box, with nothing after it.
\setbox2\vbox{\noindent A$$A$$\par}
\message{<display-last=\the\ht2>}

% A display after a box rather than after text.
\setbox3\vbox{\hbox{A}$$A$$}
\message{<display-after-box=\the\ht3>}

% `\everypar` must not fire inside the formula. If it does, the token it
% inserts is typeset in math mode and the count goes up.
\count0=0
\everypar{\global\advance\count0 by1 }
\setbox4\vbox{\noindent A$$A$$A\par}
\message{<everypar-count=\the\count0>}
\everypar{}

% The display's skips are on the box's *vertical* list, so changing them
% changes the box's height by exactly that much.
\abovedisplayskip=10pt \belowdisplayskip=11pt
\abovedisplayshortskip=10pt \belowdisplayshortskip=11pt
\setbox5\vbox{\noindent A$$A$$A\par}
\message{<with-skips=\the\ht5>}
\abovedisplayskip=0pt \belowdisplayskip=0pt
\abovedisplayshortskip=0pt \belowdisplayshortskip=0pt

\end
